% formal/capacity-derivatives.tex — derivatives and sensitivity
% Sources:
%   historical exp-sys-landscape/sensitivity-analysis/math.tex

\subsection{Derivatives of the Systolic Ratio}
\label{sec:sys-optimization}

% Relocated from historical exp-sys-landscape/sensitivity-analysis/math.tex
% (experiment deleted 2026-04-03, superseded by gradient-ascent-general/gradient-ascent-products).
% These lemmas document library/src/derivatives.rs.
% Label sec:sys-optimization kept for backward compatibility with
% cross-references in formal/hko-symmetry-gradient-structure.tex.

For a polytope $K = \{x \in \mathbb{R}^4 \mid \langle a_k, x \rangle \le 1,\; k=1,\dots,F\}$
with dual vertices $a_k \in \mathbb{R}^4 \setminus \{0\}$,
the systolic ratio decomposes as
$\mathrm{sys} = c_{\mathrm{EHZ}}^2 / (2\,\mathrm{vol})$.
The gradient $\partial/\partial a_k \in \mathbb{R}^4$ is a single vector
per facet, without the manifold constraints that arise in the unit-normal
parameterization.

\begin{theorem*}[Envelope theorem for constrained optimization]
Let $f(x;\theta)$ be a smooth objective and $g_i(x;\theta) = 0$,
$i = 1,\dots,m$, smooth equality constraints, with parameter $\theta$.
Let $x^*(\theta)$ be a non-degenerate KKT point with multipliers
$\lambda^*(\theta)$, and let
$\mathcal{L}(x,\lambda;\theta) = f(x;\theta) - \sum_i \lambda_i g_i(x;\theta)$
be the Lagrangian.
Then the optimal value $f^*(\theta) = f(x^*(\theta);\theta)$ satisfies
\[
  \frac{\partial f^*}{\partial \theta}
  = \frac{\partial \mathcal{L}}{\partial \theta}\bigg|_{x=x^*,\,\lambda=\lambda^*}
  = \frac{\partial f}{\partial \theta}\bigg|_{x^*}
    - \sum_i \lambda_i^* \frac{\partial g_i}{\partial \theta}\bigg|_{x^*}.
\]
\end{theorem*}

\paragraph{Capacity derivative.}

% [TODO: JÖRN - verify statement and proof — replaces old lem:cap-derivative + lem:cap-derivative-normal]
\begin{unverified}
\begin{lemma}[Capacity derivative w.r.t.\ dual vertices]
\label{lem:cap-derivative}
Fix an orbit $(S,\sigma)$ with $|S| = m$ facets. The orbit's action is
$A(S,\sigma;a) = 1/(2\,Q^*)$, where $Q^*$ is the maximum of the KKT problem
\[
  Q(\beta) = \tfrac{1}{2}\beta^T H \beta
  \quad\text{subject to}\quad
  N^T\beta = 0,\; \eta^T\beta = 1,
\]
with $H$ the symmetric matrix from Lemma~\ref{lem:H-quadratic},
$N = [a_{\sigma(1)} | \cdots | a_{\sigma(m)}]$, and
$\eta = \mathbf{1}$ (all ones), so the normalization constraint is
$\sum_i \beta_i = 1$.

Let $\mu \in \mathbb{R}^4$ be the Lagrange multiplier for the closing
constraint $N^T\beta = 0$ at the optimum. If facet $k$ appears in the
orbit at position $i_0$ (i.e.\ $\sigma(i_0) = k$), define
\[
  P_{i_0} = \sum_{i < i_0} \beta_i^*\,a_{\sigma(i)}.
\]
Then
\[
  \frac{\partial Q^*}{\partial a_k}
  = \beta_{i_0}^*\,\bigl[J_0(2\,P_{i_0} + \beta_{i_0}^*\,a_k) + \mu\bigr],
\]
and the capacity derivative is
\[
  \frac{\partial c}{\partial a_k}
  = -\frac{1}{2\,(Q^*)^2}\,\frac{\partial Q^*}{\partial a_k}.
\]
If $k \notin S$, then $\partial c / \partial a_k = 0$.
\end{lemma}

\begin{proof}
The Lagrangian is
$\mathcal{L} = \tfrac{1}{2}\beta^T H\beta + \mu^T N^T\beta + \xi(\mathbf{1}^T\beta - 1)$,
with stationarity condition $H\beta + N\mu + \mathbf{1}\xi = 0$.
Since $\eta = \mathbf{1}$ does not depend on~$a_k$, the dual vertices enter
through $H$ and~$N$ only.
By the envelope theorem, $\partial Q^*/\partial a_k$ involves only the
explicit dependence of $H$ and $N$ on $a_k$:
\[
  \frac{\partial Q^*}{\partial a_k}
  = \frac{1}{2}\frac{\partial((\beta^*)^T H\beta^*)}{\partial a_k}
    + \mu^T\frac{\partial(N^T\beta^*)}{\partial a_k}.
\]
For the $N$-term: only column~$i_0$ of $N$ depends on $a_k$, giving
$\mu^T(\partial N^T\beta^*/\partial a_k) = \beta_{i_0}^*\,\mu$
(as a vector in $\mathbb{R}^4$).

For the $H$-term: $Q = \frac{1}{2}\beta^T H\beta
= \sum_{j < i} \beta_i \beta_j\,\omega_0(a_{\sigma(j)}, a_{\sigma(i)})$
by Lemma~\ref{lem:H-quadratic}. Differentiating the terms involving $a_k$
($\sigma(i_0) = k$) and using the closing constraint
$\sum_i \beta_i^* a_{\sigma(i)} = 0$ to simplify yields
$\frac{1}{2}\partial(\beta^T H\beta)/\partial a_k
= \beta_{i_0}^*\,J_0(2P_{i_0} + \beta_{i_0}^* a_k)$.
(The detailed sign chain involves the antisymmetry of $\omega_0$ and the
bilinear identity $\omega_0(u,v) = (J_0 u) \cdot v$.)

Combining the H-term and N-term:
$\partial Q^*/\partial a_k
= \beta_{i_0}^*\,[J_0(2P_{i_0} + \beta_{i_0}^* a_k) + \mu]$.
Since $c = 1/(2Q^*)$, the chain rule gives
$\partial c/\partial a_k = -(\partial Q^*/\partial a_k) / (2\,(Q^*)^2)$.
\end{proof}
\end{unverified}

\paragraph{Volume derivative.}

% [TODO: JÖRN - verify statement and proof — replaces old lem:vol-derivative + lem:vol-derivative-normal]
\begin{unverified}
\begin{lemma}[Volume derivative w.r.t.\ dual vertices]
\label{lem:vol-derivative}
For a convex polytope $K = \{x : \langle a_k, x \rangle \le 1\}$ with
$0 \in \operatorname{int}(K)$,
\[
  \frac{\partial\,\mathrm{vol}(K)}{\partial a_k}
  = -\frac{S_k}{\|a_k\|^3}\,a_k
    + \frac{1}{\|a_k\|}\,\nabla_{n_k}\mathrm{vol},
\]
where $S_k = \mathrm{vol}_3(F_k)$ is the $3$-dimensional volume of facet
$F_k$, $n_k = a_k/\|a_k\|$ is the unit normal, $h_k = 1/\|a_k\|$ is the
support height, and
$\nabla_{n_k}\mathrm{vol} = -S_k\,(\bar{x}_k - h_k\,n_k)$
is the tangential volume gradient with $\bar{x}_k$ the centroid of~$F_k$.
\end{lemma}

\begin{proof}
The volume depends on $a_k$ through the height $h_k = 1/\|a_k\|$ and
the unit normal $n_k = a_k/\|a_k\|$.  By the chain rule,
\[
  \frac{\partial\,\mathrm{vol}}{\partial a_k}
  = \frac{\partial\,\mathrm{vol}}{\partial h_k}\,
    \frac{\partial h_k}{\partial a_k}
  + \left(\frac{\partial\,\mathrm{vol}}{\partial n_k}\right)^{\!T}
    \frac{\partial n_k}{\partial a_k}.
\]

\emph{Height term.}
Moving facet $k$ outward by $\delta$ enlarges $K$ by a slab of
approximate volume $\delta \cdot S_k$, so
$\partial\,\mathrm{vol}/\partial h_k = S_k$.
Since $h_k = \|a_k\|^{-1}$, we have
$\partial h_k/\partial a_k = -a_k/\|a_k\|^3$.

\emph{Normal term.}
Tilting facet $k$ by $\varepsilon\delta$ (with $\delta \perp n_k$)
replaces the constraint $\langle n_k, x \rangle \le h_k$ with
$\langle n_k + \varepsilon\delta, x \rangle \le h_k$, changing the volume by
$-\varepsilon\,S_k\,(\bar{x}_k \cdot \delta) + O(\varepsilon^2)$.
The tangential gradient (projected to $T_{n_k}S^3$) is
$\nabla_{n_k}\mathrm{vol} = -S_k\,(\bar{x}_k - h_k\,n_k)$.
Since $\partial n_k/\partial a_k = (I - n_k n_k^T)/\|a_k\|$, and
$\nabla_{n_k}\mathrm{vol} \perp n_k$, the projection acts as the identity:
$(\nabla_{n_k}\mathrm{vol})^T (I - n_k n_k^T)/\|a_k\|
= \nabla_{n_k}\mathrm{vol} / \|a_k\|$.

Combining both terms yields the stated formula.
\end{proof}
\end{unverified}

\emph{Derivative cross-checks.}
Both the capacity and volume analytical derivatives are cross-checked against
central finite differences ($\varepsilon = 10^{-6}$) on fixture polytopes
(hypercube, simplex) with relative tolerance~$10^{-4}$.

% [TODO: JÖRN - verify statement and proof sketch — was prop:capacity-piecewise-smooth]
\begin{unverified}
\begin{proposition}[Capacity as piecewise-smooth function of dual vertices]
\label{prop:capacity-piecewise-smooth}
Let $K(a)$ be a \emph{simple} polytope (every vertex determined by exactly~$4$
facets) with variable dual vertices $a = (a_1,\dots,a_F)$.
\begin{enumerate}
\item[\textup{(a)}]
  \emph{Combinatorial type stability.}
  For each $a_0$ such that $K(a_0)$ is a simple polytope, the vertex-facet
  incidence graph is locally constant: there exists $\varepsilon > 0$ such that
  $K(a)$ has the same combinatorial type for all $\|a - a_0\| < \varepsilon$.
  In particular, the set of valid orbit nodes $(S,\sigma)$ surviving
  the HK2017 pruning is constant in this neighborhood.

\item[\textup{(b)}]
  \emph{Per-orbit smoothness.}
  Fix an orbit $(S,\sigma)$ whose KKT problem has a unique non-degenerate
  maximum. Then the action $A(S,\sigma;\cdot)$ is smooth in~$a$,
  with derivatives given by Lemma~\ref{lem:cap-derivative}.

\item[\textup{(c)}]
  \emph{Capacity at generic parameters.}
  The capacity
  \[ c_{\mathrm{EHZ}}(a) = \min_{(S,\sigma)} A(S,\sigma;a) \]
  is continuous.
  At parameters where the minimizing orbit is unique, $c_{\mathrm{EHZ}}$ is
  differentiable with gradient given by that orbit's envelope theorem
  formula. The unique-minimizer case is generic: for any two orbits,
  the equation $A(S_1,\sigma_1) = A(S_2,\sigma_2)$ cuts out a hypersurface
  in the $4F$-dimensional parameter space, so the set of parameters where
  any two orbits tie has measure zero.

\item[\textup{(d)}]
  \emph{Capacity at switching boundaries.}
  When $r \ge 2$ orbits achieve minimum action simultaneously,
  the capacity is non-smooth.
  The directional derivative in direction~$d = (d_{a_1},\dots,d_{a_F})$ is
  \[
    D_d\,c_{\mathrm{EHZ}}
    = \min_{i=1,\dots,r} \nabla_a A(S_i,\sigma_i) \cdot d,
  \]
  i.e.\ the most pessimistic rate of change among the tied orbits.
\end{enumerate}
\end{proposition}

\begin{proof}[Proof sketch]
(a) Each vertex of a simple polytope is the unique solution of a $4 \times 4$
linear system $A_v\,x = \mathbf{1}$ where $A_v$ is the matrix of the four
determining dual vertices. Since $A_v$ is invertible, vertices move
continuously (indeed smoothly) in~$a$.
For each vertex, the strict inequalities
$\langle a_j, v \rangle < 1$ for non-determining facets are open conditions, so they
persist under small perturbations of~$a$.
Similarly, for ridge-adjacent facets, $\omega_0(a_i, a_j) \neq 0$ is an open
condition, so the directed adjacency graph (and hence the pruned orbit set)
is locally constant.
(b) Follows from the envelope theorem applied to the KKT problem, whose
optimum is locally unique and smooth by the implicit function theorem
(given LICQ and second-order sufficient conditions, which hold generically).
(c) Continuity: $\min$ of continuous functions.
Differentiability: when the minimizer is unique, $c_{\mathrm{EHZ}} =
A(S^*,\sigma^*;\cdot)$ locally, hence smooth by~(b).
Genericity: $A(S_1,\sigma_1;a) = A(S_2,\sigma_2;a)$ is one equation in~$4F$
unknowns, so the tie locus is a hypersurface; finitely many orbit pairs, so
the union has measure zero.
(d) Standard result for the directional derivative of a pointwise minimum
of smooth functions.
\end{proof}
\end{unverified}

% Jörn: math approved (b01bf24) — Remark
\begin{remark}[Simplicity is generic]
\label{rem:simplicity-generic}
A polytope in dual-vertex representation is non-simple when $5$ or more facet
hyperplanes $\langle a_k, x \rangle = 1$ meet at a single point, which
imposes one additional equation beyond the $4$ needed to determine a vertex.
This is a codimension-$1$ condition in the space of dual vertices, so non-simple
polytopes form a measure-zero subset: a random polytope is almost surely simple.
\end{remark}

% [TODO: JÖRN - verify statement — was cor:sys-derivative]
\begin{unverified}
\begin{corollary}[Systolic ratio derivative]
\label{cor:sys-derivative}
The chain rule applied to
$\mathrm{sys} = c_{\mathrm{EHZ}}^2 / (2\,\mathrm{vol})$ gives
\[
  \frac{\partial\,\mathrm{sys}}{\partial a_k}
  = \frac{c_{\mathrm{EHZ}}}{\mathrm{vol}}
    \cdot \frac{\partial\,c_{\mathrm{EHZ}}}{\partial a_k}
  - \frac{\mathrm{sys}}{\mathrm{vol}}
    \cdot \frac{\partial\,\mathrm{vol}}{\partial a_k},
\]
where the partial derivatives of $\mathrm{vol}$ and $c_{\mathrm{EHZ}}$ are
given by Lemmas~\ref{lem:vol-derivative} and~\ref{lem:cap-derivative}.
\end{corollary}
\end{unverified}
